\section{Scope, attribution, and limitations}\label{sec:scope}

The assertion for arbitrary ordinary six-dimensional Kakeya sets is a consequence of the explicitly stated four-dimensional finite input
and the named Katz--Tao recurrence.
For a complete verification, the reader must also examine the original proof of the four-dimensional finite interface and the advanced inputs it cites.
Author attribution, archival hashes, numerical experiments, and internal review labels are not proofs of correctness.

Three levels must be distinguished: the full-dimensional theorem for four-dimensional sets;
the four-dimensional finite directional-mass estimate for labelled original marks;
and its higher-dimensional consequences through the recurrence.
This paper uses the second level as its actual input.
Its output applies uniformly to each new finite family; the families need not arise from one common set across scales.
However, it imposes a long-mark threshold on each line and therefore does not assert a sharp Kakeya maximal estimate for arbitrarily small shading densities.

Directional selection retains the original angular upper bound as an unnormalized subfamily bound.
Passing to a normalized probability law requires division by the retained mass.
The average shading lower bound first undergoes heavy-line selection before entering the per-line input,
so\eqref{eq:finite-sharp-budget} indeed has two factors of $\lambda$.
Projected directions become separated only after selection in the collision graph;
projected Frostman control is not inherited for free.
Arbitrary fixed ambient spaces and real angular parameters are genuine requirements of the recurrence, rather than merely changes in terminology.

Relative to the ordinary six-dimensional Hausdorff baseline explicitly listed in\cite{KT,HRZ,GLX},
equation~\eqref{eq:comparison} gives a fixed positive improvement in the exponent.
The four-dimensional manuscript\cite[Section 11.2]{OAI4} explicitly records the lower bound $4$ obtained by projection in higher ambient dimensions.
The bounded comparison through the stated verification cutoff did not identify an explicit report of the same $17/4$
or a concrete result of the same type that dominates it.
Finite searches, however, establish neither absolute priority nor current optimality, and do not exclude implicit consequences elsewhere.

We do not propose a new Katz--Tao recurrence or claim independent credit for the four-dimensional manuscript or its advanced inputs.
This paper does not settle full Hausdorff dimension in five dimensions or in general dimension, nor does it construct a genuine counterexample.
The manuscript is intended for further external mathematical review.
Internal cross-review is distinguished strictly from external peer review.
